\subsection{一个自同态的相似不变量}

把 VII 第 8 页定理 1 与第 9 页命题 2 中关于有限生成挠模的分解翻译过来，便得到：

命题3. --- 设 E 是交换域 K 上有限 n 维的向量空间，u 是 E 的一个自同态；对每个首一不可约多项式 \(p(\mathbf{X})\)，设 \(M_{p}\) 为由 E 中满足 \((p(u))^{k}\) . x = 0（对某个整数 k）的元素 x 组成的向量子空间。则 \(M_{p}\) 在 u 下封闭，向量空间 E 是诸 \(M_{p}\) 的直和，并且存在多项式 \(s_{p}\)，使得对所有 \(x \in E\)，x 在 \(M_{p}\) 中的分量等于 \(s_{p}(u)\) . x。

注记1. --- 显然，u 在 \(M_{p}\) 上限制的最小多项式是整除 u 的最小多项式的 p 的最大幂。此外我们有 \(s_{p}(u).x = x\)（给定 \(x \in M_{p}\)），由此立即推出，若 \(M_{p} \neq 0\)，则 \(s_{p}\) 与 p 互素。

类似地，由 VII 第 19 页定理 2，模 \(\mathbf{E}_u\) 同构于循环模 \(\mathrm{F}_j = \mathrm{K}[\mathrm{X}] / \mathfrak{a}_j\) ( \(1 \leqslant j \leqslant r\) )的直和，其中诸理想 \(\mathfrak{a}_j\) 异于 \(\mathrm{K}[\mathrm{X}]\)，且 \(\mathfrak{a}_j \subset \mathfrak{a}_{j+1}\)；并且诸 \(\mathfrak{a}_j\) 由这些条件确定。此外，由于 \(\mathrm{E}_u\) 是挠模，我们有 \(\mathfrak{a}_1 \neq (0)\)；由于 E 的维数为 n，我们有 \(r \leqslant n\)。令 \(\mathfrak{a}_j = (h_j)\) ( \(1 \leqslant j \leqslant r\) )，其中 \(h_j\) 是首一多项式，并考虑由下式定义的多项式序列 \((q_i)\) ( \(1 \leqslant i \leqslant n\) )：

\[
\left\{ \begin{array}{l l} q _ {i} (\mathbf {X}) = 1 & \text { if } \quad i \leqslant n - r \\ q _ {i} (\mathbf {X}) = h _ {n - i + 1} (\mathbf {X}) & \text { if } \quad n - r <   i \leqslant n. \end{array} \right.\tag{3}
\]

显然，多项式 \(q_{i}\) 决定多项式 \(h_{j}\)，反之亦然；并且 \(E_{u}\) 同构于这 n 个模 \(\mathrm{K}[\mathrm{X}]/(q_{i})\) 的直和，其中前 n-r 个为 0。

换言之：

命题 4. --- 设 E 是交换域 K 上有限 n 维的向量空间，u 是 E 的一个自同态。则存在 n 个首一多项式 \(q_{i}(X) \in K[X]\) ( \(1 \leqslant i \leqslant n\) )，使得 \(q_{i}\) 整除 \(q_{i+1}\)（给定 \(1 \leqslant i \leqslant n-1\)），且 E 是 n 个子空间 \(V_{i}\) ( \(1 \leqslant i \leqslant n\) )的直和，这些子空间在 u 下封闭、（关于 u）循环，并且使得 u 在 \(V_{i}\) 上限制的最小多项式等于 \(q_{i}\) ( \(1 \leqslant i \leqslant n\) )。多项式 \(q_{i}\) 由这些条件唯一确定，且 \(q_{n}\) 是 u 的最小多项式 q。

注记 2. --- 由上述命题，存在 E 的一组基，使得 u 的矩阵 U 具有形式

\[
\left( \begin{array}{c c c c c} A _ {n - r + 1} & 0 & \dots & 0 & 0 \\ 0 & A _ {n - r + 2} & \dots & 0 & 0 \\ \cdots & \cdots & \cdots & \cdots & \cdots \\ 0 & 0 & \dots & A _ {n - 1} & 0 \\ 0 & 0 & & 0 & A _ {n} \end{array} \right)
\]

其中每个矩阵 \(A_{i}\) 具有形式 (2)（取 \(g(\mathbf{X}) = q_{i}(\mathbf{X})\) )（参见 VII，第 29-30 页）。

定义2. --- 在命题 4 的记号下，\(n\) 个首一多项式 \(q_{i}(\mathbf{X}) (1 \leqslant i \leqslant n)\) 称为自同态 \(u\) 的相似不变量。

于是，第 n 个相似不变量 \(q_{n}\) 是 u 的最小多项式（命题 4）；换言之，要使多项式 \(p(\mathbf{X}) \in \mathbf{K}[\mathbf{X}]\) 满足 \(p(u) = 0\)，必须且只需 p 是 \(q_{n}\) 的倍式。

推论1. --- 设 K 是一个域，设 E 和 E' 是 K 上的两个有限维向量空间，并设 u（相应地，u'）是 E（相应地，E'）的一个自同态。则 u 和 u' 相似（VII，第 29 页）当且仅当它们具有相同的相似不变量。事实上，u 和 u' 相似当且仅当 K{[}X{]}-模 \(E_{u}\) 与 \(E_{u'}\) 同构。

推论2. --- 设 u 是域 K 上有限维向量空间 E 的一个自同态，令 \((q_{1}, \ldots, q_{n})\) 为 u 的相似不变量族，设 L 为 K 的一个扩域，设 \(\mathrm{E}_{(\mathrm{L})} = \mathrm{L} \otimes_{\mathrm{K}} \mathrm{E}\) 为由 E 经标量扩张得到的 L-向量空间，并令 \(u_{(\mathrm{L})} = 1_{\mathrm{L}} \otimes u\) 为由 u 诱导的 \(\mathrm{E}_{(\mathrm{L})}\) 的自同态。则 \(u_{(\mathrm{L})}\) 的相似不变量是像 \(\bar{q}_{1}, \ldots, \bar{q}_{n}\)，即 \(q_{1}, \ldots, q_{n}\) 在 L{[}X{]} 中的像。

这由命题 4 以及如下事实直接得出：L{[}X{]}-模 \(\mathrm{E}_{(\mathrm{L})u_{(\mathrm{L})}}\) 与 \((\mathrm{K}[\mathrm{X}]/(q_{i}))_{(\mathrm{L})}\) 分别同构于 L{[}X{]}⊗ \(_{K[X]}\) \(E_{u}\) 与 L{[}X{]}/( \(\bar{q}_{i}\) )。

设 U 为域 K 上 n 阶方阵，其中 K 为交换域。则由 U 定义的 \(K^{n}\) 自同态的相似不变量称为 U 的相似不变量。由上述推论 1 可知，两个方阵相似当且仅当它们具有相同的相似不变量；并且若 u 是 K 上有限维向量空间 E 的自同态，U 是 u 关于 E 的某个基 B 的矩阵，则 U 与 u 的相似不变量相同。由上述推论 1 和 2，我们有：

推论3. --- 设 U 和 V 是某个交换域 K 上两个 n 阶方阵。若存在 K 的某个扩域 \(K'\) 上的可逆方阵 P 使得 \(V = P^{-1}UP\)，则存在 K 上的可逆方阵 Q 使得 \(V = Q^{-1}UQ\)。

设 E 是交换域 K 上的一个有限维向量空间，设 \((e_{i})_{1 \leq i \leq n}\) 是 E 的一组基，且设 u 是 E 的一个自同态。则由 VII，第 29 页的正合序列 (1)，与 u 相伴的 K{[}X{]}-模 \(E_{u}\) 同构于自由 K{[}X{]}-模 E{[}X{]}（以 \((1 \otimes e_{i})\) 为基）被 E{[}X{]} 在 K{[}X{]}-线性映射 X - \(\bar{u}\) 下的像所除得的商。于是 u 的相似不变量 \(q_{i}(X)\)（VII，第 32 页，定义 2）就是 X - \(\bar{u}\) 的不变因子（VII，第 22 页）。因此 VII，第 22 页的命题 6 蕴含：

命题5. --- 设 E 是交换域 K 上有限 n 维的向量空间，设 u 是 E 的一个自同态，并设 U 是 u 关于 E 的某一组基的矩阵。那么对每个满足 \(1 \leq m \leq n\) 的整数 m，乘积

\[
d _ {m} (\mathbf {X}) = q _ {1} (\mathbf {X}) q _ {2} (\mathbf {X}) \dots q _ {m} (\mathbf {X})
\]

即前 \(m\) 个相似不变量（\(u\) 的）之积，等于 \(m\) 阶子式（取自矩阵 \(XI_{n} - U\)）的最大公因数。

推论1. --- 设 u 是交换域 K 上有限 n 维向量空间的自同态，其特征多项式为 \(\chi_{u}(\mathbf{X})\)，相似不变量为 \(q_{i}(\mathbf{X})\) ( \(1 \leqslant i \leqslant n\) )。则

\[
\chi_ {u} (\mathbf {X}) = q _ {1} (\mathbf {X}) q _ {2} (\mathbf {X}) \dots q _ {n} (\mathbf {X}).
\]

推论2. --- 在推论 1 的记号中，设 \(q(\mathbf{X})\) 为 \(u\) 的最小多项式；则 \(q(\mathbf{X})\) 整除 \(\chi_u(\mathbf{X})\)，且 \(\chi_u(\mathbf{X})\) 整除 \(q(\mathbf{X})^n\)。特别地，\(u\) 的最小多项式和特征多项式具有相同的根，且这些根就是 \(u\) 的特征值。

由于 \(q(\mathbf{X}) = q_{n}(\mathbf{X})\)，显然 \(q(\mathbf{X})\) 整除 \(\chi_{u}(\mathbf{X})\)。另一方面，由于每个 \(q_{i}\) 整除 q，它们的乘积 \(\chi_{u}\) 整除 \(q^{n}\)。

推论 3. --- 一个自同态 \(u\) 是幂零的，当且仅当其特征多项式具有形式 \(X^n\)。

这由推论 2 立即得出。

现在让我们改写 VII，第 24 页的命题 9，该命题给出模分解为不可分解子模的直和。

命题6. --- 设 E 是交换域 K 上有限 n 维的向量空间，且 u 是 E 的一个自同态。则 E 是诸子空间 \(E_{k}\) 的直和，这些子空间在 u 下封闭且关于 u 循环，使得 u 在 \(E_{k}\) 上限制的最小多项式具有形式 \(p_{k}^{n(k)}\)，其中 \(p_{k}\) 是不可约多项式，且 \(E_{k}\) 不能表示为两个各自在 u 下封闭的非零子空间的直和。对每个首一不可约多项式 \(p \in K[X]\) 及每个整数 \(n \geqslant 1\)，数 \(m(p^{n})\)（即这样的子空间 \(E_{k}\) 的个数：在任何这样的分解中，\(p^{n}\) 是 u 在 \(E_{k}\) 上限制的最小多项式）是唯一确定的。

诸 \(p_{k}^{n(k)}\) 确定 u 的相似不变量，反之亦然；我们可以借助 VII，第 25 页，注记 2 和 3 中解释的过程从一个得到另一个。此外，我们可以立即从命题 6 中所考虑的分解得到命题 3 和 4 中所考虑的分解。

注意，首一不可约多项式 \(p \in K[X]\)（它使得 \(m(p^{n}) > 0\) 对某个整数 \(n \geqslant 1\) 成立）恰好是 u 的最小多项式的首一不可约因子。因此，与相似不变量不同，这些多项式一般依赖于我们工作于其中的域 K。

